\bibliographystyle{plain}
\begin{thebibliography}{99}\small

\bibitem{anac2011}
T.~Baarslag, K.~Fujita, E.~H. Gerding, K.~Hindriks, T.~Ito, N.~R. Jennings,
C.~Jonker, S.~Kraus, R.~Lin, V.~Robu, and C.~R. Williams.
Evaluating practical negotiating agents: results and analysis of the 2011
international competition.
\emph{Artificial Intelligence}, 198:73--103, 2013.

\bibitem{baarslag2013acceptance}
T.~Baarslag, K.~Hindriks, and C.~Jonker. Acceptance conditions in automated
negotiation. In \emph{Complex Automated Negotiations: Theories, Models, and
Software Competitions}, volume 435 of \emph{Studies in Computational
Intelligence}, pages 95--111. Springer, 2013.

\bibitem{baarslag2014decoupling}
T.~Baarslag, K.~Hindriks, M.~Hendrikx, A.~Dirkzwager, and C.~Jonker.
Decoupling negotiating agents to explore the space of negotiation strategies.
In \emph{Novel Insights in Agent-based Complex Automated Negotiation}, volume
535 of \emph{Studies in Computational Intelligence}, pages 61--83. Springer,
2014.

\bibitem{baarslag2016survey}
T.~Baarslag, M.~J.~C. Hendrikx, K.~V. Hindriks, and C.~M. Jonker.
Learning about the opponent in automated bilateral negotiation: a comprehensive
survey of opponent modeling techniques.
\emph{Autonomous Agents and Multi-Agent Systems}, 30(5):849--898, 2016.

\bibitem{crawford1982strategic}
V.~P. Crawford and J.~Sobel. Strategic information transmission.
\emph{Econometrica}, 50(6):1431--1451, 1982.

\bibitem{kamenica2011bayesian}
E.~Kamenica and M.~Gentzkow. Bayesian persuasion.
\emph{American Economic Review}, 101(6):2590--2615, 2011.

\bibitem{myerson1983}
R.~B. Myerson and M.~A. Satterthwaite. Efficient mechanisms for bilateral trading.
\emph{Journal of Economic Theory}, 29(2):265--281, 1983.

\bibitem{ochsroth1989}
J.~Ochs and A.~E. Roth. An experimental study of sequential bargaining.
\emph{American Economic Review}, 79(3):355--384, 1989.

\bibitem{rubinstein1982}
A.~Rubinstein. Perfect equilibrium in a bargaining model.
\emph{Econometrica}, 50(1):97--109, 1982.

\bibitem{shapira2024glee}
E.~Shapira, O.~Madmon, M.~Tennenholtz, and R.~Reichart.
GLEE: A unified framework and benchmark for language-based economic environments.
\emph{arXiv:2410.05254}, 2024. \url{https://arxiv.org/abs/2410.05254}

\end{thebibliography}

% =====================================================================
%  SUPPLEMENTARY MATERIAL -- not part of the 4-page main body.
% =====================================================================
\appendix
\newpage
\section{Supplementary figure and tables}

% Appendix floats are numbered A1, A2, ... so a reference from the main text
% reads unambiguously as supplementary material.
\setcounter{table}{0}
\setcounter{figure}{0}
\renewcommand{\thetable}{A\arabic{table}}
\renewcommand{\thefigure}{A\arabic{figure}}

\begin{table}[H]
\centering\small
\caption{Seller-script ladder (persuasion buyer seat): forward hit rate on the
\emph{next} purchase, conditioned on what that exact message template has
already delivered earlier in the same game (6{,}185 buyer replays). The 5.3\%
row sits below the break-even bar $1/m$ in every cell, which is what makes the
template veto strictly dominated.}
\label{tab:ladder}
\begin{tabular}{@{}lrlr@{}}
\toprule
Prior record of this template & hit rate & Prior record & hit rate \\
\midrule
no prior           & 70.4\% & all-low $\times1$ & 33.8\% \\
all-high $\times1$ & 81.3\% & all-low $\times2$ & 19.3\% \\
all-high $\times2$ & 84.8\% & all-low $\times3$ & \tbf{5.3\%} \\
all-high $\times3$ & 91.2\% & mixed             & 63.7\% \\
\bottomrule
\end{tabular}
\end{table}

\begin{figure}[H]
\centering
\begin{tikzpicture}[x=1cm,y=1cm,line width=0.4pt]
\draw[gridline] (0.000,0.344) -- (12.672,0.344);
\node[axlab,anchor=east] at (-0.080,0.344) {2200};
\draw[gridline] (0.000,0.918) -- (12.672,0.918);
\node[axlab,anchor=east] at (-0.080,0.918) {2250};
\draw[gridline] (0.000,1.492) -- (12.672,1.492);
\node[axlab,anchor=east] at (-0.080,1.492) {2300};
\draw[gridline] (0.000,2.066) -- (12.672,2.066);
\node[axlab,anchor=east] at (-0.080,2.066) {2350};
\node[axlab,anchor=north] at (0.000,-0.043) {0h};
\node[axlab,anchor=north] at (2.880,-0.043) {2h};
\node[axlab,anchor=north] at (5.760,-0.043) {4h};
\node[axlab,anchor=north] at (8.640,-0.043) {6h};
\node[axlab,anchor=north] at (11.520,-0.043) {8h};
\draw[axline] (0.000,0.000) -- (12.672,0.000);
\draw[axline] (0.000,0.000) -- (0.000,2.410);
\fill[band] (0.000,1.064) -- (0.248,1.324) -- (0.505,1.221) -- (0.801,1.258) -- (1.098,1.333) -- (1.393,1.407) -- (1.659,1.527) -- (1.955,1.443) -- (2.251,1.430) -- (2.547,1.538) -- (2.842,1.470) -- (3.109,1.479) -- (3.374,1.457) -- (3.638,1.570) -- (3.933,1.596) -- (4.229,1.556) -- (4.493,1.565) -- (4.757,1.852) -- (5.055,1.735) -- (5.320,1.805) -- (5.616,1.752) -- (5.912,1.837) -- (6.209,2.232) -- (6.474,2.073) -- (6.771,2.045) -- (7.066,2.009) -- (7.363,2.020) -- (7.659,1.987) -- (7.925,2.047) -- (8.222,2.029) -- (8.487,2.101) -- (8.752,2.141) -- (9.018,2.074) -- (9.313,2.155) -- (9.579,2.346) -- (9.874,2.359) -- (10.170,2.286) -- (10.465,2.050) -- (10.729,2.024) -- (10.993,2.104) -- (11.289,2.099) -- (11.584,2.029) -- (11.848,2.022) -- (12.144,1.937) -- (12.440,2.000) -- (12.440,1.475) -- (12.144,1.369) -- (11.848,1.509) -- (11.584,1.565) -- (11.289,1.546) -- (10.993,1.622) -- (10.729,1.880) -- (10.465,1.818) -- (10.170,1.667) -- (9.874,1.723) -- (9.579,1.844) -- (9.313,1.853) -- (9.018,1.900) -- (8.752,1.752) -- (8.487,1.527) -- (8.222,1.550) -- (7.925,1.626) -- (7.659,1.620) -- (7.363,1.754) -- (7.066,1.561) -- (6.771,1.546) -- (6.474,1.373) -- (6.209,1.403) -- (5.912,1.430) -- (5.616,1.161) -- (5.320,1.136) -- (5.055,1.041) -- (4.757,0.990) -- (4.493,0.969) -- (4.229,1.042) -- (3.933,1.147) -- (3.638,1.073) -- (3.374,1.222) -- (3.109,1.081) -- (2.842,0.991) -- (2.547,0.951) -- (2.251,0.941) -- (1.955,0.773) -- (1.659,0.683) -- (1.393,0.572) -- (1.098,0.586) -- (0.801,0.463) -- (0.505,0.439) -- (0.248,0.315) -- (0.000,0.084) -- cycle;
\draw[c1] (0.000,0.388) -- (0.248,0.315) -- (0.505,0.674) -- (0.801,0.598) -- (1.098,0.645) -- (1.393,0.730) -- (1.659,0.899) -- (1.955,1.033) -- (2.251,1.002) -- (2.547,1.088) -- (2.842,1.349) -- (3.109,1.344) -- (3.374,1.323) -- (3.638,1.250) -- (3.933,1.434) -- (4.229,1.321) -- (4.493,1.357) -- (4.757,1.431) -- (5.055,1.264) -- (5.320,1.324) -- (5.616,1.161) -- (5.912,1.445) -- (6.209,1.403) -- (6.474,1.373) -- (6.771,1.546) -- (7.066,1.742) -- (7.363,1.849) -- (7.659,1.652) -- (7.925,1.708) -- (8.222,1.771) -- (8.487,1.661) -- (8.752,1.930) -- (9.018,1.947) -- (9.313,1.853) -- (9.579,1.844) -- (9.874,1.927) -- (10.170,1.954) -- (10.465,2.050) -- (10.729,1.880) -- (10.993,1.622) -- (11.289,1.546) -- (11.584,1.565) -- (11.848,1.509) -- (12.144,1.369) -- (12.440,1.475);
\node[ptlab,c1,anchor=west] at (12.490,1.475) {B$_1$};
\draw[c2] (0.000,0.396) -- (0.248,0.778) -- (0.505,0.709) -- (0.801,0.825) -- (1.098,0.742) -- (1.393,0.811) -- (1.659,0.867) -- (1.955,0.773) -- (2.251,1.040) -- (2.547,0.951) -- (2.842,1.181) -- (3.109,1.081) -- (3.374,1.222) -- (3.638,1.073) -- (3.933,1.147) -- (4.229,1.042) -- (4.493,0.969) -- (4.757,0.990) -- (5.055,1.041) -- (5.320,1.136) -- (5.616,1.189) -- (5.912,1.551) -- (6.209,1.610) -- (6.474,1.742) -- (6.771,1.817) -- (7.066,2.009) -- (7.363,2.020) -- (7.659,1.971) -- (7.925,2.013) -- (8.222,1.931) -- (8.487,1.888) -- (8.752,2.141) -- (9.018,2.054) -- (9.313,1.870) -- (9.579,1.874) -- (9.874,1.723) -- (10.170,1.667) -- (10.465,1.818) -- (10.729,1.935) -- (10.993,1.883) -- (11.289,1.757) -- (11.584,1.782) -- (11.848,1.675) -- (12.144,1.739) -- (12.440,1.513);
\node[ptlab,c2,anchor=west] at (12.490,1.513) {B$_2$};
\draw[c3] (0.000,1.064) -- (0.248,1.324) -- (0.505,1.221) -- (0.801,1.258) -- (1.098,1.333) -- (1.393,1.407) -- (1.659,1.527) -- (1.955,1.443) -- (2.251,1.430) -- (2.547,1.538) -- (2.842,1.470) -- (3.109,1.479) -- (3.374,1.457) -- (3.638,1.418) -- (3.933,1.287) -- (4.229,1.283) -- (4.493,1.143) -- (4.757,1.380) -- (5.055,1.397) -- (5.320,1.418) -- (5.616,1.391) -- (5.912,1.430) -- (6.209,1.720) -- (6.474,1.724) -- (6.771,1.705) -- (7.066,1.561) -- (7.363,1.805) -- (7.659,1.754) -- (7.925,1.626) -- (8.222,1.550) -- (8.487,1.527) -- (8.752,1.752) -- (9.018,1.900) -- (9.313,2.145) -- (9.579,2.346) -- (9.874,2.359) -- (10.170,2.286) -- (10.465,2.019) -- (10.729,1.931) -- (10.993,1.857) -- (11.289,1.995) -- (11.584,1.935) -- (11.848,1.720) -- (12.144,1.744) -- (12.440,1.509);
\node[ptlab,c3,anchor=west] at (12.490,1.509) {B$_3$};
\draw[c4] (0.000,0.345) -- (0.248,0.626) -- (0.505,0.653) -- (0.801,0.636) -- (1.098,0.591) -- (1.393,0.572) -- (1.659,0.683) -- (1.955,0.944) -- (2.251,1.137) -- (2.547,1.271) -- (2.842,1.406) -- (3.109,1.430) -- (3.374,1.426) -- (3.638,1.570) -- (3.933,1.596) -- (4.229,1.556) -- (4.493,1.565) -- (4.757,1.852) -- (5.055,1.735) -- (5.320,1.581) -- (5.616,1.589) -- (5.912,1.533) -- (6.209,1.629) -- (6.474,1.684) -- (6.771,1.570) -- (7.066,1.718) -- (7.363,1.754) -- (7.659,1.620) -- (7.925,1.657) -- (8.222,1.722) -- (8.487,1.838) -- (8.752,1.984) -- (9.018,2.031) -- (9.313,2.155) -- (9.579,2.155) -- (9.874,2.190) -- (10.170,2.069) -- (10.465,1.954) -- (10.729,1.881) -- (10.993,2.104) -- (11.289,2.099) -- (11.584,2.029) -- (11.848,1.913) -- (12.144,1.686) -- (12.440,1.733);
\node[ptlab,c4,anchor=west] at (12.490,1.733) {B$_4$};
\draw[c5] (0.000,0.084) -- (0.248,0.482) -- (0.505,0.439) -- (0.801,0.463) -- (1.098,0.586) -- (1.393,0.702) -- (1.659,0.878) -- (1.955,1.036) -- (2.251,0.941) -- (2.547,1.062) -- (2.842,0.991) -- (3.109,1.184) -- (3.374,1.272) -- (3.638,1.313) -- (3.933,1.262) -- (4.229,1.438) -- (4.493,1.430) -- (4.757,1.365) -- (5.055,1.723) -- (5.320,1.805) -- (5.616,1.752) -- (5.912,1.837) -- (6.209,2.232) -- (6.474,2.073) -- (6.771,2.045) -- (7.066,1.859) -- (7.363,1.985) -- (7.659,1.987) -- (7.925,2.047) -- (8.222,2.029) -- (8.487,2.101) -- (8.752,2.131) -- (9.018,2.074) -- (9.313,2.147) -- (9.579,2.275) -- (9.874,1.945) -- (10.170,1.912) -- (10.465,1.975) -- (10.729,2.024) -- (10.993,1.907) -- (11.289,1.991) -- (11.584,1.958) -- (11.848,2.022) -- (12.144,1.937) -- (12.440,2.000);
\node[ptlab,c5,anchor=west] at (12.490,2.000) {B$_5$};
\draw[spreadarrow] (6.209,1.403) -- (6.209,2.232);
\node[note,anchor=south west] at (6.309,2.011) {\textbf{72 rating points apart,}};
\node[note,anchor=north west] at (6.309,2.011) {byte-identical code};
\end{tikzpicture}

\caption{\tbf{Five byte-identical builds, 8.6\,h, 45 simultaneous polls,
${\approx}3{,}000$ games each.} Band = min--max envelope per poll (mean spread
48.1, max 88.0). Per family the same builds spread 136.8 / 161.8 / 109.8 points
on average and 292.9 at peak, while their \emph{behaviour} differed by
${\le}0.010$ of the pot.}
\label{fig:noise}
\end{figure}

\begin{table}[H]
\centering\small
\caption{Fleet volume, 12--30 August 2026, recovered from the local turn logs.
Only the account's best slot is ranked, which is what makes the other four
available as simultaneous controls.}
\label{tab:volume}
\begin{tabular}{@{}lrlr@{}}
\toprule
Family & unique games & Fleet total & \\
\midrule
Bargaining  & 89{,}641 & logged turns         & 3{,}114{,}356 \\
Negotiation & 76{,}490 & unique games         & 264{,}432 \\
Persuasion  & 98{,}301 & reconstructable runs & 49 \\
\bottomrule
\end{tabular}
\end{table}

\begin{table}[H]
\centering\small
\caption{Break-even recovery (technique C) over 7{,}104 known-horizon and
31{,}208 open-horizon bargaining games. The whole bargaining deficit is one
cell, and it is the accept side: replaying the shipped decision rule over
26{,}773 logged decision points (98.79\% control fidelity) shows the stonewall
term setting the bar as low as 0.118, so it signs 0.30--0.50 offers in the
opening rounds. \tbf{Break-even} is the outcome at which mean $\Delta R$ crosses
zero, which is the outcome ranking at my \emph{own} standing $R_t$; it equals the
field median only where $R_t{\approx}2000$, which holds in these cells but should
not be read as a field median in general.}
\label{tab:breakeven}
\begin{tabular}{@{}llrrrl@{}}
\toprule
Horizon & Seat & I bank & Break-even & Gap & Diagnosis \\
\midrule
open     & $p_1$ & 0.530 & 0.476 & $+0.054$ & healthy \\
open     & $p_2$ & 0.501 & 0.463 & $+0.038$ & healthy \\
known 12 & $p_1$ & 0.435 & \tbf{0.816} & $\mathbf{-0.381}$ & accept bar collapses to 0.118 \\
known 12 & $p_2$ & 0.687 & 0.449 & $+0.239$ & healthy \\
\bottomrule
\end{tabular}
\end{table}

\begin{table}[H]
\centering\small
\caption{Open-horizon bargaining: the break-even outcome tracks Rubinstein SPE,
and I sit under it in every player-1 cell (technique C, 38{,}312 games). The
exception is $\delta_{\text{us}}{=}\delta_{\text{them}}{=}0.95$, where SPE is
0.513 but break-even is only 0.475, so demanding SPE there buys no-deals, which
is why the shipped floor targets the measured value rather than the equilibrium.
As in Table~\ref{tab:breakeven}, break-even reads as a field median only where
$R_t{\approx}2000$.}
\label{tab:spe}
\begin{tabular}{@{}rrrrrr@{}}
\toprule
$\delta_{\text{us}}$ & $\delta_{\text{them}}$ & mine & break-even & SPE & gap \\
\midrule
0.95 & 0.80 & 0.712 & 0.825 & 0.833 & $-0.11$ \\
0.95 & 0.90 & 0.536 & 0.625 & 0.690 & $-0.09$ \\
0.90 & 0.80 & 0.595 & 0.675 & 0.714 & $-0.08$ \\
0.90 & 0.90 & 0.449 & 0.525 & 0.526 & $-0.08$ \\
1.00 & 0.90 & 0.692 & 0.775 & 1.000 & $-0.08$ \\
0.80 & 0.80 & 0.477 & 0.525 & 0.556 & $-0.05$ \\
\bottomrule
\end{tabular}
\end{table}

\begin{table}[H]
\centering\small
\caption{Five of the six changes that replicated, the sixth omitted for space.
Each removed a \emph{strictly dominated} choice, an option plainly worse than an
available alternative with no trade-off to price. Nothing else did.}
\label{tab:ledger}
\begin{tabular}{@{}p{0.26\linewidth}p{0.70\linewidth}@{}}
\toprule
Change & Evidence \\
\midrule
Complete-info ultimatum & I was asking 0.650 of a \emph{known} surplus on the one round they cannot refuse; 0.99 was signable. $+3.42\sigma$, $n{=}214$ \\
Rung-aware pricing & Price at the recovered valuation rungs, not guesses. $+3.9/{+}4.1\sigma$, replicated \\
Free-clock accept floor & At $\delta{=}1$ refusing is costless. ${+}0.0458$ of pot, $n{=}108$, $7.7\sigma$, none worse \\
Open-horizon cap collapse & The server zeroes both sides and never says so. Caught 11 of 11 \$0 games \\
Negative-signal parser & 270 buys made against messages that read as refusals; \tbf{95.6\% low quality} against a 27.5\% base rate \\
\bottomrule
\end{tabular}
\end{table}

\begin{table}[H]
\centering\small
\caption{Bounded complete-information negotiation: the same seat under two
horizons, over 10{,}740 games. Grinding to the last round is where the
family's losses are, and deal rate cannot see it, because the round-10 dump
\emph{is} a deal.}
\label{tab:endgame}
\begin{tabular}{@{}lrlr@{}}
\toprule
Cell / outcome & $n$ & banked / buyer value & rating\,/\,game \\
\midrule
$\text{max\_rounds}{=}1$, seller, complete info  & 488 & asks 0.991, signed 99.2\% & $\mathbf{+6.10}$ \\
$\text{max\_rounds}{=}10$, seller, complete info & 469 & (split below) & $\mathbf{-4.63}$ \\
\quad closed by round 8  & 66  & 0.2365 & $+5.9$ \\
\quad ground to round 10 & 403 & \tbf{0.0433} & $\mathbf{-6.4}$ \\
\bottomrule
\end{tabular}
\end{table}

\section{Extended limitations}
\label{app:lim}

Three limitations were compressed out of \S8. \tbf{(i) The persuasion buyer seat
admits no per-round counterfactual even in principle}: quality on declined rounds
is never revealed. Technique (D) recovers it by subtraction instead. With a known
prior $p$ of high quality, a bought fraction $b$ of rounds and a realised hit rate
$h$ on those, the quality rate on the rounds I passed is estimated by
$\widehat{q}_{\text{declined}}{=}(p-bh)/(1-b)$, since the two strata must
reconstitute the prior in expectation. That answers only the aggregate question ``am I leaking in this cell'',
never ``should I have bought round~7''.
\tbf{(ii) The offline simulator is optimistic.} I built a local bargaining
engine (real rules, my seat played by the shipped code, the other seat by a
field fitted to 15{,}788 real games plus eight adversarial archetypes). It
reproduces per-cell median settlement rounds but scores ${\approx}{+}0.03$ of pot
high, with a ${\pm}45$-rating noise floor per world at $n{=}4{,}000$, so I used
it only to choose which arm deserved a live window and never shipped on it
alone. \tbf{(iii) Operational failure modes cost me more than any strategy
error.} Three recur: a deploy that writes arm flags back into the repository
working tree (hit five times, and the reason 27 tests fail on my final-night
tree); an entry point outside the copied package, so sibling agents silently ran
stale code; and force-killing an agent mid-game, which abandons in-flight games
at the 5th percentile and can trip a 30-minute queue ban. Draining with a
deadline rather than a kill is the only safe stop.

\tbf{(iv) Opponent modelling is stable but unreachable where it would pay.}
\S8 reports split-half $r{=}0.786$ across the 19 opponent\,$\times$\,cell
profiles with ${\ge}25$ purchases in both halves, which is real but is also the
whole usable sample. Three things block the rest. Identity is withheld in
49.9\% of games; the named half spreads 16{,}498 games over 227 opponents and
15 cells, about five games per profile, and behaviour differs by cell so they
cannot be pooled; and quality is redacted on every round I passed (0 of 42{,}616
unbought rounds ever showed it), so a reputation table has \emph{zero} entries
in precisely the cells where I buy nothing and most need one. That is why the
shipped buyer reads the seller's rationing rate within the current game, which
is visible whether or not I buy, instead of a cross-game profile.

\section{A note on the Simpson's paradox I nearly shipped on}
\label{app:simpson}

Four analyses failed this way. Worked through here is the one that came closest
to shipping; the others were a closed-game deal rate, a buyer-quality rate, and
an abandonment cost, each of which conditioned its denominator on the very
outcome it was trying to measure.

A late candidate rule, ``take a recommendation only from a seller who has
already spent a sale to decline one'', looked refuted when pooled across all
fifteen $(m,p)$ cells: subsequent dud rates of 0.285 / 0.287 / 0.286 / 0.315
for sellers who had declined 0 / 1 / 2 / 3--4 of their first four rounds. Split
by cell it reverses and becomes monotone: inside $(m{=}1.20, p{=}0.80)$,
$n{=}9{,}932$, the same buckets read 0.160 / 0.102 / 0.081 / \tbf{0.069}
against a break-even of $1-1/m = 0.167$. The pooled number is dominated by easy
cells where everyone buys everything regardless. I did not ship it, for the
reason the rest of this paper argues: at 2.2\% of my games it was worth
${\approx}{+}2.6$ rating points, far below what the remaining window could
measure, and a fourth simultaneous change would have destroyed my ability to
read the three already running.

\section{Three properties of the scoring rule}
\label{app:scoring}

\tbf{The aggregate rewards balance, and I optimised the wrong object for two
weeks.} Overall rank is the \emph{unweighted mean} of the three family ratings,
and only an account's best slot is ranked. Judged per family, negotiation was my
best return per game; judged per \emph{overall} point it was nearly worthless,
because ${+}300$ more negotiation is ${+}100$ overall and hard, while ${+}180$ in
each weak family is ${+}120$ and was what the arms already targeted. At the time
I noticed, I sat 11th at 2301.5 with negotiation 2875 (4th best on the board) and
bargaining 2021; the 4th-placed account held that same bargaining rating.
\tbf{Nobody in the top ten was balanced}, which is the strongest structural
incentive the scoring rule creates and the one I was slowest to see.

\tbf{The reference pool defends itself against sabotage.} A natural worry is
that spare agent slots could be spent depressing rivals' ratings. Besides being
prohibited and trivially detectable, it is self-defeating on the mechanism:
percentile is scored against payoffs earned on the same configuration, and
competition games join that pool as they are played, so every no-deal forced on
an opponent injects a bottom-of-distribution payoff that \emph{raises} the
percentile of every normal payoff scored afterwards. Attacking the field inflates
it. There is also no opponent selection anywhere in the API, so the damage would
spray uniformly across the 227 opponents I faced.

\tbf{The underlying data will come from the organisers rather than from me.}
Every number here is computed from my own turn logs, which also carry the
transcripts of the agents I was matched against, and those are not released
alongside the code. The organisers have said they will publish the full game
record for every agent, which supersedes anything I could have attached: it
covers the whole field rather than one account, and it comes from the party that
ran the games. The analysis scripts are released now, so the method can be
audited ahead of the data it will run on.

\tbf{The garden of forking paths is wide here.} One persuasion arm, analysed
three defensible ways on the same data, read $-3.11\sigma$ across 3 cells,
$-0.82\sigma$ across 33, and $+3.02\sigma$ across 44, depending only on the
inclusion threshold for how many games a configuration needed to enter the
sample. Pre-registering the cell filter before looking is not optional at this
noise level.

\end{document}
